\documentclass[10pt,journal,compsoc]{IEEEtran}

% --- 宏包加载 ---
\usepackage[utf8]{inputenc}
\usepackage[T1]{fontenc}
\usepackage{amsmath, amssymb, amsfonts, amsthm}
\usepackage{bm}
\usepackage{booktabs}
\usepackage{algorithm}
\usepackage{algorithmic}
\usepackage{graphicx}
\usepackage{cite}
\usepackage{hyperref}

% --- 自定义命令 ---
\newcommand{\SynOps}{\text{SynOps}}

\begin{document}


\section{Notation and SNN Forward Model (Event-Based Detection)}
Let $\mathcal{D} = \{(x_n, y_n)\}_{n=1}^N$ be an event-based detection dataset, where $x_n$ is an event stream and $y_n$ is the set of labeled boxes and classes[cite: 2]. We train an SNN detector $f_{\Theta}$ (baseline anchored on EAS-SNN) that maps events to detection predictions $\hat{y} = f_{\Theta}(x)$[cite: 3].

\subsection{Discrete-time LIF Dynamics with Soft Reset}
For layer $l$, neuron $i$, and time step $t \in \{1, \dots, T\}$, denote input spikes $s^{(l-1)}[t] \in \{0,1\}^{N_{l-1}}$, membrane potential $u^{(l)}[t] \in \mathbb{R}^{N_l}$, and output spikes $s^{(l)}[t] \in \{0,1\}^{N_l}$[cite: 4]. A standard discrete LIF update is defined as[cite: 5]:
\begin{equation}
u^{(l)}[t] = \beta^{(l)} u^{(l)}[t-1] + I^{(l)}[t] - \theta^{(l)} s^{(l)}[t-1],
\end{equation}
\begin{equation}
s^{(l)}[t] = H\left(u^{(l)}[t] - \theta^{(l)}\right),
\end{equation}
where $\beta^{(l)} \in (0,1)$ is the leak factor, $\theta^{(l)}$ is the firing threshold, and $H(\cdot)$ is the Heaviside step function[cite: 5]. The $-\theta s[t-1]$ term implements \textbf{soft reset}, ensuring robustness under low-precision state representations[cite: 5]. For a convolutional synapse[cite: 6]:
\begin{equation}
I^{(l)}[t] = W^{(l)} * s^{(l-1)}[t],
\end{equation}
with $W^{(l)} \in \mathbb{R}^{C_{\text{out}}^{(l)} \times C_{\text{in}}^{(l)} \times K \times K}$[cite: 6].

\subsection{Training with Surrogate Gradients}
To handle the non-differentiability of $H(\cdot)$, we employ surrogate gradients during training[cite: 7]:
\begin{equation}
\frac{\partial s^{(l)}[t]}{\partial u^{(l)}[t]} \approx \phi\left(u^{(l)}[t] - \theta^{(l)}\right),
\end{equation}
where $\phi(\cdot)$ is a smooth proxy derivative[cite: 7]. This enables end-to-end optimization for event-based backbones[cite: 8].

\subsection{Detection Objective}
We optimize a differentiable detection loss $\mathcal{L}_{\text{det}}$ as a tractable handle for mAP[cite: 9, 11]:
\begin{equation}
\mathcal{L}_{\text{det}}(\Theta) = \mathcal{L}_{\text{cls}}(\hat{p}, p) + \lambda_{\text{box}}\mathcal{L}_{\ell_1}(\hat{b}, b) + \lambda_{\text{iou}}\mathcal{L}_{\text{iou}}(\hat{b}, b),
\end{equation}
where $\hat{p}$ and $\hat{b}$ represent class probabilities and bounding box coordinates, respectively[cite: 10].

\section{Hardware-Aware NAS Under FPGA Constraints}
NAS must optimize accuracy under discrete hardware constraints, as DSP and BRAM resources are width-aligned[cite: 12, 13].

\subsection{Design Variables}
The deployable design is jointly described by:
\begin{itemize}
    \item \textbf{Architecture ($a$)}: Per-layer structural choices (channels, kernels, depth)[cite: 14].
    \item \textbf{Structured mask ($m$)}: SIMD-aligned block activation variables[cite: 15].
    \item \textbf{Precision ($q = (b_W, b_U)$)}: Per-layer bitwidths for weights and membrane states[cite: 16].
\end{itemize}

\subsection{Hardware-Observable Indicators}
On UltraScale+ FPGAs, the key observables driving latency and energy are DSP48E2 slice count ($R_{\text{DSP}}$), BRAM36 block count ($R_{\text{BRAM}}$) with width granularity, and LUT usage ($R_{\text{LUT}}$)[cite: 17, 18, 19]. These determine the achieved clock frequency $f_{\text{clk}}$ and memory bandwidth $B_{\text{mem}}$[cite: 22, 23].

\subsection{Spike-Rate--Aware Compute}
Let the average spike rate be $\rho^{(l)}$[cite: 25]. The effective synaptic operations are[cite: 27]:
\begin{equation}
\SynOps^{(l)}(a,m) \approx T \cdot \rho^{(l-1)} \cdot \SynOps^{(l)}_{\text{dense}} \cdot \kappa^{(l)}(m),
\end{equation}
where $\kappa^{(l)}(m)$ captures the fraction of blocks kept after structured pruning[cite: 27].

\subsection{Discrete Resource Models}
Required DSPs for layer $l$ are modeled as[cite: 28]:
\begin{equation}
R_{\text{DSP}}^{(l)}(a,m,q) = \left\lceil \frac{P^{(l)}_{\text{req}}(a,m)}{\nu_{\text{DSP}}(b_W^{(l)})} \right\rceil,
\end{equation}
where $\nu_{\text{DSP}}$ is the SIMD packing factor[cite: 28]. BRAM usage considering width penalty $\psi(w)$ is:
\begin{equation}
R_{\text{BRAM}}(w,d) = \left\lceil \frac{w \cdot d}{C_{\text{BRAM}}} \right\rceil \cdot \psi(w).
\end{equation}
Total BRAM usage is the sum of weight storage $R_{\text{BRAM},W}$ and membrane state storage $R_{\text{BRAM},U}$[cite: 33, 34].

\subsection{Constrained NAS Objective}
We optimize the architecture through Lagrangian penalties[cite: 40]:
\begin{equation}
\begin{aligned}
\min_{a,m,q,\Theta} \quad & \mathcal{L}_{\text{det}}(\Theta; a,m,q) + \lambda_{\text{lat}}[\text{Lat} - \tau_{\text{lat}}]_{+} \\
& + \lambda_{\text{eng}}[\text{Eng} - \tau_{\text{eng}}]_{+} + \sum_{r} \lambda_r[R_r - \tau_r]_{+},
\end{aligned}
\end{equation}
where $[\cdot]_{+} = \max(0, \cdot)$ and $\tau$ denotes platform budgets[cite: 40].

\section{Block-Wise Structured Pruning}
To avoid the "unstructured sparsity trap," we align pruning with SIMD lanes[cite: 44, 45].

\subsection{SIMD-Aligned Grouping}
Weights are partitioned into blocks $g \in \mathcal{G}^{(l)}$ matching the hardware vectorization width $B_{\text{simd}}$[cite: 46]. The pruned weights are $\widetilde{W}^{(l)} = W^{(l)} \odot M^{(l)}$[cite: 47].

\subsection{Cost-Aware Block Selection (Knapsack)}
Block retention is formulated as a multi-constraint knapsack problem:
\begin{equation}
\max_{m_g \in \{0,1\}} \sum_g v_g m_g \quad \text{s.t.} \quad \sum_g c_{g,r} m_g \leq \tau_r,
\end{equation}
where $v_g$ is the Taylor importance and $c_g$ is the discrete resource increment[cite: 54, 56].

\section{Quantization for SNNs}
We quantize both weights and membrane potentials to address the "membrane storage wall"[cite: 63, 73].

\subsection{Uniform Quantization}
A symmetric uniform quantizer is defined as[cite: 65]:
\begin{equation}
Q_b(x;\Delta) = \Delta \cdot \text{clip}\left(\left\lfloor \frac{x}{\Delta} \right\rceil, -2^{b-1}, 2^{b-1}-1\right).
\end{equation}

\subsection{State-Aware Shift-Friendly Leakage}
To maximize efficiency, we constrain leak factors to $\beta^{(l)} = 1 - 2^{-n^{(l)}}$[cite: 71]. The fully quantized recurrence becomes:
\begin{equation}
u_q^{(l)}[t] = Q_{b_U} \left( u_q[t-1] - (u_q[t-1] \gg n) + I_q - \theta_q s[t-1] \right).
\end{equation}

\section{Unified HAPQ Procedure}
The pipeline follows: 1) Train baseline, 2) Hardware calibration, 3) Constrained NAS, 4) Structured pruning, and 5) QAT for weights and membrane states[cite: 80, 81, 82, 83].

\end{document}